%% 自動生成（gen.py）
\bktocchap{3-1}{\R{変}{か}わる\R{量}{りょう}と\R{変}{か}わらない\R{量}{りょう}}{ずっと\R{変}{か}わらない\R{数}{かず}と、\R{問題}{もん|だい}ごとに\R{決}{き}まる\R{数}{かず}}{7}{1}{1}{24}{C}
\bktocchap{3-2}{\R{場所}{ば|しょ}を\R{数}{かず}で\R{表}{あらわ}す：\R{座標}{ざ|ひょう}}{「どこ」を2つの\R{数}{かず}で\R{伝}{つた}えよう}{6}{1}{2}{48}{C}
\bktocchap{3-3}{\R{関数}{かん|すう}は「\R{入}{い}れると\R{出}{で}てくる」ルール}{\R{表}{ひょう}・グラフ・\R{式}{しき}は、\R{同}{おな}じことを\R{言}{い}っている}{5}{1}{2}{76}{C}
\bktocchap{3-4}{グラフをかこう・\R{読}{よ}もう}{\R{実験}{じっ|けん}の\R{点}{てん}に、なめらかな\R{線}{せん}を\R{引}{ひ}こう}{7}{1}{1}{104}{C}
\bktocchap{3-5}{グラフの\R{傾}{かたむ}きは「\R{変化}{へん|か}の\R{速}{はや}さ」}{\R{坂}{さか}が\R{急}{きゅう}なほど、\R{速}{はや}く\R{変}{か}わっている}{1}{2}{2}{128}{C}
\bktocchap{3-6}{2\R{本}{ほん}の\R{直線}{ちょく|せん}が\R{交}{まじ}わるところ：\R{連立方程式}{れん|りつ|ほう|てい|しき}}{ふたりが\R{出会}{で|あ}うのは、グラフの\R{交}{まじ}わる\R{点}{てん}}{5}{2}{2}{162}{C}
